\chapter{Software and Hardware Requirement Specification}
For the successful implementation of an AI-enabled education system, it is important to develop a solid technical base which balances out the need for high-performance training and low-weight production. This chapter discusses the technical requirements needed for supporting the multi-agent system and the LLM-based pipeline.

\section{Functional Requirements}
In order to meet the goals relating to SDG 4, 9, and 10, the system will adhere to the following key functionalities:
\begin{itemize}
	\item \textbf{Real-time Tutoring:} The system will provide instant responses to students' queries by deploying agents that utilize large language models (LLMs).
	\item \textbf{Personalized Learning Paths:} The multi-agent architecture will evaluate student learning outcomes and modify the complexity of the learning content accordingly.
	\item \textbf{User Profiling:} The system will keep track of the activities of users in every session, ensuring that there is always a current profile. 
	\item \textbf{Information Retrieval:} he system will provide a function that allows for searching for and retrieving relevant information using vector-based search queries.
\end{itemize}

\section{Non-functional Requirements}
\begin{itemize}
	\item \textbf{Scalability:} The software architecture must be able to handle an increase in the number of simultaneous users without sacrificing performance.
	\item \textbf{Latency:} For an efficient learning experience, the agent latency must, wherever possible, remain below two seconds.
	\item \textbf{Availability:} The production environment will be accessible through Internet connectivity and have limited downtime.
	\item \textbf{Cost Efficiency:} The development process will employ open-source architectures and cloud computing systems that can be supported by free-tier accounts.
\end{itemize}

\section{Hardware Requirements}
Hardware requirements are categorized into two different stages: The demanding training stage, and the well-optimized production stage.

\subsection{For Training and Development}
\begin{itemize}
	\item \textbf{Processor:} Intel Core i5 or i7 series (8th gen or above) or AMD Ryzen 5 or 7 series.
	\item \textbf{RAM:} A minimum requirement of 8 GB RAM is needed; 16 GB RAM is highly recommended to facilitate the processing of complex natural language processing models.
	\item \textbf{Graphics Card:} An NVIDIA graphics card of GTX 1650 or higher version is recommended, which supports CUDA operations to enhance the efficiency of model training, or else Google Colab GPU of T4/L4 version can be used.
	\item \textbf{Disk Storage:} The disk storage should be at least 200 GB.
	\item \textbf{Internet Connection:} A fast and stable internet connection is necessary to download pre-trained weights of model files.
\end{itemize}

\subsection{Production and Deployment}
\begin{itemize}
	\item \textbf{Internet:} A reliable internet connection is necessary to maintain seamless API requests and users' usage.
	\item \textbf{Disk Storage:} Disk storage needs to be at least 10 GB.
\end{itemize}

\section{Software Requirements}
The software architecture revolves around the AI software stack in the Python ecosystem and modern microservices software stack for the backend.

\subsection{Development and Machine Learning}
\begin{itemize}
	\item \textbf{Programming Language:} Python 3.12 or above for machine learning projects; Java and Python for backend.
	\item \textbf{IDE:} Visual Studio Code, Jupyter Notebook, Google Colab.
	\item \textbf{Machine Learning / Deep Learning Libraries:} TensorFlow, scikit-learn for modeling.
	\item \textbf{NLP Libraries:} Huggingface, SpaCy, and NLTK for natural language processing.
	\item \textbf{Preprocessing / Data Libraries:} Pandas, NumPy.
\end{itemize}

\subsection{System Design and Orchestration}
\begin{itemize}
	\item \textbf{Multi-Agent Systems:} Consider using LangChain, CrewAI, or AutoGen.
	\item \textbf{Database:} Adopt either NoSQL (MongoDB) or Relational (PostgreSQL) database systems alongside vector databases like FAISS or Qdrant for Retrieval-Augmented Generation (RAG).
	\item \textbf{Backend:} Develop high-performing API endpoints using FastAPI or Quarkus.
	\item \textbf{Deployment:} Create Docker containers and deploy the working model using Hugging Face Spaces.
\end{itemize}

\section{Cost Estimation}
The economic approach taken in this project focuses on maximizing efficiency and sustainability by leveraging modern open-source technologies to ensure that initial investments are kept to a minimum. Expected costs include the following:
\begin{enumerate}
	\item \textbf{Compute:} Code development is done using institution-based systems, as well as Google Colab, which has a free tier option.
	\item \textbf{Hosting:} Hosting is kept light by utilizing the free tier of Hugging Face Spaces.
	\item \textbf{Software:} ll software requirements are covered by open-source software (Python, FastAPI, LangChain).
\end{enumerate}
\textbf{Estimated Total Cost:} Minimal, if any. Costs will be incurred for computing power in case of an increase in users beyond what is available in the free tier.